\appendix{Фрагменты исходного кода программы}

В настоящем приложении приведены фрагменты исходного кода ключевых модулей серверной части разработанной системы (репозиторий \texttt{ispgestorserver}). Полный исходный код серверной части размещен в указанном репозитории, клиентской части~--~в репозитории \texttt{ispgestoradmin}.

Фрагмент сервиса интеграции с маршрутизатором MikroTik (\texttt{app/Services/MikroTikService.php}):
\begin{lstlisting}[language=PHP, firstnumber=1]
<?php

namespace App\Services;

use RouterOS\Client;
use RouterOS\Query;
use Exception;
use Illuminate\Support\Facades\Log;

class MikroTikService
{
    protected $client;

    public function __construct(?Client $client)
    {
        $this->client = $client;
    }

    public function runQuery(Query $query): array
    {
        if (!$this->client) {
            return [];
        }
        return $this->client->query($query)->read();
    }

    // Agregar una IP a una address-list (p. ej. la lista 'morosos')
    public function addIpToAddressList(string $ip, string $listName, string $comment = ''): array
    {
        if (!$this->client) {
            return ['success' => false, 'message' => 'Cliente MikroTik no inicializado'];
        }

        try {
            // Verificar si ya existe
            $print = new Query('/ip/firewall/address-list/print');
            $print->where('list', $listName);
            $print->where('address', $ip);
            $exists = $this->client->query($print)->read();

            if (!empty($exists)) {
                return ['success' => true, 'message' => 'La IP ya está en la lista', 'already_exists' => true];
            }

            $add = new Query('/ip/firewall/address-list/add');
            $add->equal('list', $listName);
            $add->equal('address', $ip);
            if ($comment) {
                $add->equal('comment', $comment);
            }

            $result = $this->client->query($add)->read();

            if (isset($result['!trap'])) {
                return ['success' => false, 'message' => $result['!trap'][0]['message'] ?? 'Error desconocido de MikroTik'];
            }

            return ['success' => true, 'message' => 'IP agregada a la lista correctamente'];
        } catch (Exception $e) {
            Log::error('MikroTik Add Address List Error: ' . $e->getMessage());
            return ['success' => false, 'message' => $e->getMessage()];
        }
    }
}
\end{lstlisting}

Фрагмент построения имени простой очереди абонента (\texttt{app/Services/MikroTikQueueSyncService.php}):
\begin{lstlisting}[language=PHP, firstnumber=1]
<?php

namespace App\Services;

use App\Models\Client;
use Illuminate\Support\Facades\Log;

class MikroTikQueueSyncService
{
    // Normaliza el nombre del abonado y le añade el documento de identidad
    protected function normalizeClientQueueName(Client $client): string
    {
        $name = trim((string) ($client->full_name ?: 'cliente_' . $client->id));
        $name = strtolower($name);
        $name = preg_replace('/\s+/', '_', $name);
        $name = preg_replace('/_+/', '_', $name);
        $name = trim($name, '_');
        $documentId = trim((string) ($client->document_id ?? ''));
        if ($documentId !== '') {
            return $name . '_' . $documentId;
        }
        return $name;
    }

    // Nombre final de la cola respetando el límite de RouterOS (64 caracteres)
    protected function buildClientQueueName(Client $client): string
    {
        $raw = $this->normalizeClientQueueName($client);
        $sanitized = strtolower($this->normalizeName($raw));
        $documentId = trim((string) ($client->document_id ?? ''));

        if ($documentId !== '' && str_ends_with($sanitized, '_' . $documentId)) {
            $max = 64;
            if (strlen($sanitized) <= $max) {
                return $sanitized;
            }

            $suffix = '_' . $documentId;
            $allowedBaseLen = max(1, $max - strlen($suffix));
            $base = substr($sanitized, 0, $allowedBaseLen);
            $base = rtrim($base, '_');
            $final = $base . $suffix;

            Log::warning('Client queue name truncated to fit RouterOS limits', [
                'client_id' => $client->id,
                'document_id' => $documentId,
                'final' => $final,
                'max_len' => $max,
            ]);

            return $final;
        }

        return $sanitized;
    }
}
\end{lstlisting}

Фрагмент сервиса приостановки и восстановления обслуживания (\texttt{app/Services/ClientSuspensionService.php}):
\begin{lstlisting}[language=PHP, firstnumber=1]
<?php

namespace App\Services;

use App\Models\Audit;
use App\Models\Client;
use App\Models\ClientPlan;
use Illuminate\Support\Facades\DB;
use Illuminate\Support\Facades\Log;

class ClientSuspensionService
{
    public function __construct(
        private readonly MikroTikService $mikrotik,
        private readonly ClientWhitelistService $whitelist,
    ) {}

    // Reactivar el servicio de un cliente tras regularizar la deuda
    public function reactivateClient(Client $client, string $reason): array
    {
        if (in_array(strtoupper($client->service_status), ['ACTIVE', 'ACTIVO'])) {
            return ['success' => true, 'already_active' => true];
        }

        $mkResult = ['skipped' => 'no_ip'];

        // Quitar la IP de la lista 'morosos' en el router
        if ($client->ip) {
            try {
                $mkResult = $this->mikrotik->removeIpFromAddressList($client->ip, 'morosos');
            } catch (\Throwable $e) {
                Log::error("ClientSuspensionService: Excepción MikroTik cliente {$client->id}: " . $e->getMessage());
                $mkResult = ['success' => false, 'error' => $e->getMessage()];
            }
        }

        // El cambio de estado en BD se aplica dentro de una transacción
        DB::transaction(function () use ($client, $reason, $mkResult) {
            $oldStatus = $client->service_status;

            $client->service_status = 'active';
            $client->save();

            ClientPlan::where('client_id', $client->id)
                ->where('status', 'suspended')
                ->update(['status' => 'active']);

            Audit::create([
                'table_name' => 'clients',
                'operation'  => 'REACTIVATE_AUTO_OP',
                'record_id'  => (string) $client->id,
                'old_values' => ['service_status' => $oldStatus],
                'new_values' => [
                    'service_status'  => 'active',
                    'reason'          => $reason,
                    'mikrotik_list'   => 'morosos',
                    'mikrotik_result' => $mkResult,
                    'timestamp'       => now()->toIso8601String(),
                ],
                'user_id'    => null,
                'ip_address' => '127.0.0.1',
            ]);
        });

        Log::info("ClientSuspensionService: Cliente {$client->id} reactivado. Razón: {$reason}");

        return ['success' => true, 'mikrotik' => $mkResult];
    }
}
\end{lstlisting}

Фрагмент сервиса автоматического биллинга (\texttt{app/Services/AutoBillingService.php}):
\begin{lstlisting}[language=PHP, firstnumber=1]
<?php

namespace App\Services;

use App\Models\Invoice;
use Illuminate\Support\Facades\DB;

class AutoBillingService
{
    // Ciclo completo: genera facturas y procesa los cobros pendientes
    public function processCompleteBilling()
    {
        $results = [];

        $results['invoices_generated'] = $this->generateMonthlyInvoices();
        $paymentRun = $this->processAutoPayments();
        $results['payments_processed'] = $paymentRun['processed'] ?? [];
        $results['summary'] = $this->getBillingSummary();

        return $results;
    }

    // Cobra una factura desde la wallet del abonado en una transacción
    public function processInvoicePayment(Invoice $invoice)
    {
        return DB::transaction(function () use ($invoice) {
            $client = $invoice->client;
            $estadoAnterior = $invoice->status;

            // Asegurar que el cliente tenga una wallet
            if (!$client->wallet) {
                $client->wallet()->create(['balance' => 0]);
                $client->load('wallet');
            }

            $wallet = $client->wallet;

            // Verificar saldo suficiente antes de cobrar
            if ($wallet->balance < $invoice->total_amount) {
                if ($invoice->status !== Invoice::STATUS_FAILED) {
                    $invoice->update(['status' => Invoice::STATUS_FAILED]);
                }

                return [
                    'success' => false,
                    'error' => 'Saldo insuficiente',
                    'previous_status' => $estadoAnterior,
                ];
            }

            try {
                $wallet->makePayment(
                    $invoice->total_amount,
                    "Pago automático - {$invoice->invoice_number}",
                    "Pago de factura {$invoice->invoice_number}"
                );

                $transaction = $wallet->transactions()->latest()->first();

                // Marcar la factura como pagada
                $invoice->markAsPaid(Invoice::PAYMENT_WALLET, $transaction->reference);

                return [
                    'success' => true,
                    'message' => 'Pago procesado exitosamente',
                    'transaction_reference' => $transaction->reference,
                    'paid_at' => now()->toDateTimeString(),
                    'previous_status' => $estadoAnterior,
                ];
            } catch (\Exception $e) {
                $invoice->update(['status' => Invoice::STATUS_FAILED]);

                return [
                    'success' => false,
                    'error' => $e->getMessage(),
                    'previous_status' => $estadoAnterior,
                ];
            }
        });
    }
}
\end{lstlisting}

Фрагмент промежуточного обработчика проверки разрешений (\texttt{app/Http/Middleware/CheckPermission.php}):
\begin{lstlisting}[language=PHP, firstnumber=1]
<?php

namespace App\Http\Middleware;

use Closure;
use Illuminate\Http\Request;
use Symfony\Component\HttpFoundation\Response;

class CheckPermission
{
    public function handle(Request $request, Closure $next, string $slug): Response
    {
        $employee = $request->user();

        if (!$employee) {
            return response()->json(['message' => 'No autenticado'], 401);
        }

        // Cargar rol y permisos si aún no están cargados
        if (!$employee->relationLoaded('role')) {
            $employee->load('role.permissions');
        } elseif ($employee->role && !$employee->role->relationLoaded('permissions')) {
            $employee->role->load('permissions');
        }

        // El rol super_admin omite toda comprobación de permisos
        if ($employee->hasRole('super_admin')) {
            return $next($request);
        }

        // El permiso 'acceso_total' también omite las comprobaciones específicas
        if ($employee->hasPermission('acceso_total')) {
            return $next($request);
        }

        // Comprobar el permiso específico requerido por la ruta
        if (!$employee->hasPermission($slug)) {
            return response()->json([
                'message' => 'No tienes permiso para realizar esta acción.',
                'required' => $slug,
            ], 403);
        }

        return $next($request);
    }
}
\end{lstlisting}

\ifВКР{
\newpage
\addcontentsline{toc}{section}{На отдельных листах (CD-RW в прикрепленном конверте)}
\noindent
\begin{tabular}{p{5.8cm}C{4.8cm}C{4.8cm}}
   Автор ВКР & \lhrulefill{\fill} & \fillcenter\Автор \\
            \setarstrut{\footnotesize}
           & \footnotesize{(подпись, дата)} & \\
            \restorearstrut
   Руководитель ВКР & \lhrulefill{\fill} & \fillcenter\Руководитель \\
            \setarstrut{\footnotesize}
           & \footnotesize{(подпись, дата)} & \\
            \restorearstrut
   Нормоконтроль & \lhrulefill{\fill} & \fillcenter\Нормоконтроль \\
            \setarstrut{\footnotesize}
           & \footnotesize{(подпись, дата)} & \\
            \restorearstrut
\end{tabular}
\vskip 2cm
\begin{center}
\textbf{Место для диска}
\end{center}
}\fi
